% ==================================================================================== %
% Example assignment template for physics/astrophysics/astronomy
% MIT (c) 2026 Taotao Yang
% ==================================================================================== %
\documentclass[11pt,a4paper]{article}
% ==================================================================================== %
% An assignment template for physics/astronomy/astrophysics
% MIT (c) 2026 Taotao Yang
% ==================================================================================== %
\usepackage[margin=1in]{geometry}
\usepackage[shortlabels]{enumitem}
\usepackage[nointegrals]{wasysym}
\usepackage[mathlines]{lineno}
\usepackage[bottom]{footmisc}
\usepackage[utf8]{inputenc}
\usepackage{mathtools}
\usepackage{listings}
\usepackage{titlesec}
\usepackage{pgfplots}
\usepackage{amsfonts}
\usepackage{graphicx}
\usepackage{fancyhdr}
\usepackage{dcolumn}
\usepackage{wrapfig}
\usepackage{physics}
\usepackage{amsmath}
\usepackage{amssymb}
\usepackage{amsthm}
\usepackage{natbib}
\usepackage{xcolor}
\usepackage{xfrac}
\usepackage{float}
\usepackage{calc}
\usepackage{bm}
% ==================================================================================== %
\definecolor{nord-bg}{HTML}{FFFFFF}
\definecolor{nord-fg}{HTML}{2E3440}
\definecolor{nord-comment}{HTML}{5E81AC}
\definecolor{nord-keyword}{HTML}{BF616A}
\definecolor{nord-string}{HTML}{A3BE8C}
\definecolor{nord-number}{HTML}{B48EAD}
\definecolor{nord-operator}{HTML}{D08770}
% ==================================================================================== %
\usepackage[colorlinks=true,urlcolor=nord-fg,linkcolor=nord-fg,citecolor=nord-fg]{hyperref}
\usepackage{cleveref}
% ==================================================================================== %
\def\[{\begin{equation}}
\def\]{\end{equation}}
\pagestyle{fancy}
\pgfplotsset{compat=1.18}
\setlength{\headheight}{14.5pt}
\renewcommand{\headrulewidth}{0.3pt}
% \newcommand{\sectionbreak}{\clearpage}
\theoremstyle{definition}
\newtheorem{problem}{Problem}
\renewcommand*{\proofname}{Solution}
% ==================================================================================== %
\titleformat{\section}{\Large\bfseries}{\thesection}{1em}{}[{\titlerule[0.0pt]}]
% ==================================================================================== %
\lstdefinestyle{nord_light}{
  backgroundcolor=\color{nord-bg},
  basicstyle=\color{nord-fg}\footnotesize\ttfamily,
  commentstyle=\color{nord-comment},
  keywordstyle=\color{nord-keyword}\bfseries,
  stringstyle=\color{nord-string},
  numberstyle=\color{nord-number},
  otherkeywords={self,cls,@classmethod,@property},
  keywordstyle={[2]\color{nord-operator}},
  frame=leftline,
  showstringspaces=false,
  breaklines=true,
  tabsize=4,
  extendedchars=true,
  numbers=left,
  numberstyle=\tiny\color{nord-fg},
  stepnumber=1,
  numbersep=10pt,
}
\lstset{style=nord_light}
% ==================================================================================== %
\bibliographystyle{apsrev}
\setcitestyle{numbers,square}
% ==================================================================================== %
% ==================================================================================== %
\fancyhf{}
\fancyhead[L]{Very Difficult Course}
\fancyhead[C]{}
\fancyhead[R]{Assignment 42}
\fancyfoot[C]{\thepage}
% ==================================================================================== %
\begin{document}
    \thispagestyle{empty}
    \begin{center}
        {\bf \Large The assignment: this looks like a title}\\
        \vspace{-0.2cm}
    \end{center}
    \begin{flushleft}
        \makebox[\linewidth]{\rule{\textwidth}{0.8pt}}
        \vspace{-0.2cm}\\
        \begin{tabular}{p{1.6cm} p{7cm} p{3cm}}
            {\bf Topic:} & Very Difficult Course (VDC) & \\
            {\bf Due:} & Jan. 01, 1970 by 23:59 hours & \\
            {\bf Author:} & James Holden &
            \href{mailto:holden@roci.edu}{holden@roci.edu}\\
            {\bf ~} & Naomi Nagata &
            \href{mailto:nagata@roci.edu}{nagata@roci.edu}\\
            {\bf ~} & Amos Burton &
            \href{mailto:burton@roci.edu}{burton@roci.edu}\\
            {\bf ~} & Alex Kamal &
            \href{mailto:kamal@roci.edu}{kamal@roci.edu}\\
            {\bf Lecturer:} & Dr. Praxideke Meng &
            \href{mailto:pmeng@ganymede.org}{pmeng@ganymede.org}\\
            {\bf ~} & Chrisjen Avasarala &
            \href{mailto:avasarala@unn.gov}{avasarala@unn.gov}\\
            {\bf Dated:} & Jan. 01, 2000 &
        \end{tabular}
        \vspace{0.2cm}\\
        \makebox[\linewidth]{\rule{\textwidth}{0.8pt}}
    \end{flushleft} 

    % \vspace{-1em}
% ==================================================================================== %
    \tableofcontents
    \clearpage
    \setcounter{page}{1} 
% ==================================================================================== %
\section{A section}
\label{sec:a_section}
    This is a section.
    
    \subsection{A subsection}
    \label{subsec:a_subsection}
    This is a subsection
% ==================================================================================== %
\section{Math}
\label{sec:math}
    Write numbered equations with labelling quickly with:
    \begin{lstlisting}[language=tex]
    \[ 
        E = m c^2
        \label{eq:einstein}
    \]
    \end{lstlisting}
    This becomes:
    \[ 
        E = m c^2
        \label{eq:einstein}
    \]
    Inline math is simply \verb|$ a =1 $|, rendered as $a=1$. 
    And we can refer to the labeled entries via \verb|\cref{}|. Let us look at \cref{eq:einstein}.
% ==================================================================================== %
\section{Figures and Tables}
\label{sec:figs_and_tabs}

    We include an image placed here:
    \begin{lstlisting}[language=tex]
    \begin{figure}[h]
        \centering
        \includegraphics[width=0.25\linewidth]{fig/test.png}
        \caption[A short caption]{A long caption with more information}
        \label{fig:an_img}
    \end{figure}
    \end{lstlisting}
    \begin{figure}[h]
        \centering
        \includegraphics[width=0.25\linewidth]{fig/test.png}
        \caption[A short caption]{A long caption with more information}
        \label{fig:an_img}
    \end{figure}

    How about a table? 
    \begin{lstlisting}[language=tex]
    begin{table}[h]
        \centering
        \begin{tabular}{l c}
            \hline
            Parameter & Requirement \\
            \hline
            entry & 1\\
            \hline
            another entry & 2\\
            \hline
        \end{tabular}
        \caption[A short caption]{A long caption}
        \label{tab:a_tab}
    \end{table}
    \end{lstlisting}
    
    \begin{table}[h]
        \centering
        \begin{tabular}{l c}
            \hline
            Parameter & Requirement \\
            \hline
            entry & 1\\
            \hline
            another entry & 2\\
            \hline
        \end{tabular}
        \caption[A short caption]{A long caption}
        \label{tab:a_tab}
    \end{table}
    
% ==================================================================================== %
\appendix
\section{Code Listing}
    \label{app: code}
    A code file can also be listed.
    \begin{lstlisting}[language=tex]
    \lstinputlisting[
        language=Python, 
        caption={code/example.py}, 
        label={list:example.py}]
    {code/example.py}
    \end{lstlisting}
    \lstinputlisting[
        language=Python, 
        caption={code/example.py}, 
        label={list:example.py}]
    {code/example.py}

\end{document}
